\documentclass[letterpaper]{article} % DO NOT CHANGE THIS
\usepackage[submission]{aaai2027}  % DO NOT CHANGE THIS
% The serif, sans-serif, and monospaced fonts are loaded automatically by
% aaai2027.sty (newtxtext, helvet, courier). DO NOT add \usepackage{times},
% \usepackage{helvet}, \usepackage{courier}, or any other font package.
\usepackage[hyphens]{url}  % DO NOT CHANGE THIS
\usepackage{graphicx} % DO NOT CHANGE THIS
\urlstyle{rm} % DO NOT CHANGE THIS
\def\UrlFont{\rm}  % DO NOT CHANGE THIS
\usepackage{natbib}  % DO NOT CHANGE THIS AND DO NOT ADD ANY OPTIONS TO IT
\usepackage{caption} % DO NOT CHANGE THIS AND DO NOT ADD ANY OPTIONS TO IT
\frenchspacing  % DO NOT CHANGE THIS
%
% These are recommended to typeset algorithms but not required.
\usepackage{algorithm}
\usepackage{algorithmic}
%
% Math support (standard, permitted).
\usepackage{amsmath}
\usepackage{amssymb}
%
% Recommended for better-looking tables
\usepackage{booktabs}
%
% Keep the \pdfinfo as shown here.
\pdfinfo{
/TemplateVersion (2027.1)
}

\setcounter{secnumdepth}{0} %May be changed to 1 or 2 if section numbers are desired.

% Title -- mixed (Title) case.
\title{DISTIL: Demonstration Distillation for Sample-Efficient Imitation Learning}

\author{
    Anonymous Submission
}
\affiliations{
    %
}

\begin{document}

\maketitle

\begin{abstract}
Expert demonstrations are the dominant cost of imitation learning, yet
interactive methods spend them reflexively: a per-state gate fires on the first
uncertain step of a rollout and the expert corrects that miss, whatever it is.
Under a fixed budget, what matters is \emph{which} failures get corrected and
\emph{how} each corrective demonstration is placed.
We present \textbf{DISTIL}, an interactive imitation-learning loop that, out
of the pool of all possible corrective demonstrations, distills the few most
informative ones, the purest corrections, and spends the fixed budget only
on those.
DISTIL flags uncertain rollouts by the learner's own per-step training loss,
describes each flagged failure with a vision language model, performs
root-cause reasoning grounded in a
task knowledge graph, clusters failures into recurring modes with a
cross-round memory, and lets a language model prescribe the round's one
corrective demonstration (a targeted correction of a recorded failure or a
bridging demonstration covering several related ones) behind a feasibility
check with re-prescription. The learner is any policy
$f_\theta$ with a per-step loss: we instantiate pure CNN and MLP classifiers on
a GridWorld task and state- and image-based diffusion policies on four robot
manipulation tasks. Under an identical 20-demonstration budget against six
interactive baselines, DISTIL attains the highest mean final success rate in
all ten task--modality settings (e.g., $96.1\%$ vs.\ $90.7\%$ for Diff-DAgger on
state Push-T; $95.3\%$ vs.\ $89.6\%$ on image Wipe), collects demonstrations
with the highest per-demonstration information gain in every setting, and
reports prescription confidences that correlate with realized policy
improvement ($r=0.82$--$0.89$).
\end{abstract}

\section{Introduction}

Behavior cloning regresses an expert's actions on the states the expert
visits~\citep{pomerleau1988alvinn,pomerleau1991efficient}, and its well-known
failure mode is covariate shift: the learner's own actions steer it off the
demonstrated support, where errors compound over the
horizon~\citep{ross2011dagger}. Interactive imitation learning fixes the
distribution by aggregating expert corrections on states the \emph{learner}
visits~\citep{ross2011dagger,celemin2022iil}, but the fix is paid for in
expert effort: every correction is a demonstration an expert must produce. The query-efficient DAgger family therefore
gates the expert behind a per-state predicate: SafeDAgger on action
discrepancy~\citep{zhang2017safedagger}, DropoutDAgger and EnsembleDAgger on
predictive uncertainty~\citep{menda2017dropoutdagger,menda2019ensembledagger},
ThriftyDAgger on novelty and risk~\citep{hoque2021thriftydagger}, and
Diff-DAgger on the diffusion policy's own training
loss~\citep{lee2025diffdagger}. All of these decide \emph{when} to interrupt a
rollout; the expert then corrects whatever failure happened to trip the gate.

\begin{figure}[t]
\centering
\includegraphics[width=0.45\columnwidth]{figures/teaser.png}
\caption{\textbf{DISTIL in a nutshell.} An LLM diagnoses the learner's
failures and \emph{prescribes} the one corrective demonstration the expert
should provide; the demonstration is added and the policy retrained. Out of
all possible corrective demonstrations, DISTIL distills the few most
informative ones and spends the fixed budget only on those.}
\label{fig:teaser}
\end{figure}

This is wasteful under the constraint that binds in practice: a \emph{fixed
demonstration budget}. When only twenty corrective demonstrations can be
collected, the first uncertain step of the current rollout is rarely the
best place to spend one. Failures repeat: a when-to-query gate can correct
the same recurring mode again and again while rarer modes starve. Placement matters too: one demonstration started \emph{between}
several related failures can cover all of them, while an on-policy correction
covers only its own rollout. Under a fixed budget the questions are
\emph{which} failures get corrected and \emph{how} each demonstration is
placed.

\textbf{DISTIL} (demonstration distillation) answers both questions with one
loop (Figures~\ref{fig:teaser} and~\ref{fig:arch}). Rollouts of the current
policy are monitored with the learner's own per-step training loss, and
uncertain episodes are flagged. A vision language model describes each
flagged failure from its start, highest-loss ($t^\star$), and end frames, and
a reasoning model assigns a root cause grounded in a task knowledge
graph. The failures are
clustered into $k$ recurring modes; a cluster memory carried across
rounds rotates attention away from modes already corrected. A language model
then \emph{prescribes} the round's single corrective demonstration: a
targeted on-policy correction of one recorded failure, or a bridging
demonstration placed at a new start state covering several related
failures. A feasibility check with re-prescription and a
deterministic fallback vets every prescription, so no budget is ever wasted. The expert supplies
exactly the prescribed demonstration, it is added to the dataset, and the
policy is updated; \emph{which} mode to correct and \emph{where} to start
the demonstration are the two arms of the hybrid prescription, committed
each round from the clustered evidence. The name is the mechanism: out of
the pool of every corrective demonstration the expert could provide, DISTIL
distills the few most informative ones, the purest corrections, and the
fixed budget is spent only on those.

DISTIL is policy-agnostic. It needs only a learner $f_\theta$ with a per-step
training loss $\ell_t$ that can be evaluated along rollouts: the cross-entropy
loss of a pure CNN (image) or MLP (state) classifier on GridWorld, and the
denoising loss of state- and image-based diffusion
policies~\citep{chi2023diffusionpolicy} on the robot tasks; the same loop
runs unchanged over both families.

We evaluate on five tasks (GridWorld $5{\times}5$, Push-T, Lift, Wipe, and
Door), each in a state and an image modality, against six interactive
baselines under an identical budget of 20 expert demonstrations. DISTIL
attains the highest mean final held-out success rate in all ten
task--modality settings, its demonstrations carry the highest
per-demonstration information gain in every setting, and its prescription
confidence predicts the realized change in success rate with Pearson
$r=0.82$--$0.89$.

\paragraph{Contributions.}
\textbf{(1)}~We reframe budget-constrained interactive imitation learning as
demonstration distillation: decide \emph{which} failure modes to correct and
\emph{how} to place each corrective demonstration, rather than only
\emph{when} to interrupt a rollout.
\textbf{(2)}~DISTIL, a policy-agnostic loop combining uncertainty flagging,
VLM failure perception, knowledge-grounded root-cause reasoning,
failure-mode clustering with cross-round cluster memory, and LLM
demonstration prescription with a feasibility-checked re-prescription loop.
\textbf{(3)}~A controlled study on five tasks $\times$ two modalities against
six baselines: DISTIL has the best mean in all ten settings at the
20-demonstration budget, the highest per-demonstration information gain everywhere,
and exposes a calibrated confidence signal ($r=0.82$--$0.89$).
\textbf{(4)}~A bridge between online and offline imitation learning:
DISTIL's feasibility-checked prescriptions are standing artifacts, so the
loop can pause after its on-policy analysis and the expert can record the
prescribed demonstrations in one batch days later; expert time is decoupled
from policy training time.

\section{Related Work}
\label{sec:related}

\subsection{Interactive Imitation Learning}
Behavior cloning dates to ALVINN~\citep{pomerleau1988alvinn,
pomerleau1991efficient}. DAgger aggregates expert labels on
learner-visited states as no-regret online
learning~\citep{ross2011dagger}; AggreVaTe and Deeply AggreVaTeD extend the
reduction to cost-to-go~\citep{ross2014aggrevate,sun2017aggrevated};
\citet{celemin2022iil} survey the field. Because vanilla DAgger queries
constantly, query-efficient variants gate \emph{when} to defer: SafeDAgger
learns a discrepancy classifier~\citep{zhang2017safedagger}; DropoutDAgger and
EnsembleDAgger threshold dropout and ensemble
uncertainty~\citep{menda2017dropoutdagger,menda2019ensembledagger,gal2016dropout,
lakshminarayanan2017ensembles}; HG-DAgger and LazyDAgger manage human-gated
control switching~\citep{kelly2019hgdagger,hoque2021lazydagger};
ThriftyDAgger budgets interventions by novelty and
risk~\citep{hoque2021thriftydagger}; DART injects noise into the expert's rollouts so offline
collection covers shifted states~\citep{laskey2017dart}; IWR and Sirius reweight samples
around human takeovers~\citep{mandlekar2020iwr,liu2023sirius}. Diffusion
policies~\citep{chi2023diffusionpolicy,ho2020ddpm,janner2022diffuser} capture
the multimodal actions that make agreement-based uncertainty brittle, and
Diff-DAgger queries on the diffusion training loss
itself~\citep{lee2025diffdagger}. All decide when to interrupt; none reasons
over the \emph{set} of a policy's failures, remembers what was corrected, or
chooses where a demonstration should begin. That gap is DISTIL's subject.

\subsection{Foundation Models for Failure Reasoning}
Language and vision-language models increasingly steer robot data collection
and control: SayCan grounds plans in affordances~\citep{ahn2022saycan}, Inner
Monologue closes a textual feedback loop~\citep{huang2022innermonologue}, and
PaLM-E and RT-2 fuse perception with language~\citep{driess2023palme,
brohan2023rt2}. A second line emits executable structure: policy
code~\citep{liang2023codeaspolicies,singh2023progprompt}, 3D value
maps~\citep{huang2023voxposer}, or reward
functions~\citep{ma2024eureka}. Closest to us is
failure reasoning: REFLECT narrates and corrects robot
failures~\citep{liu2023reflect}, AHA reasons over manipulation
failures~\citep{duan2025aha}, and Reflexion and Self-Refine iterate on model
outputs~\citep{shinn2023reflexion,madaan2023selfrefine}. These systems reason
about one trajectory at a time; DISTIL aggregates VLM-perceived failures
across rollouts, clusters them into modes, and narrows the LLM's output to
one auditable artifact: a feasibility-checked prescription for a single
corrective demonstration.

\subsection{Active Selection, Representations, and Resets}
Deciding which failures deserve a demonstration is batch active learning over
failures: pool-based selection under a budget~\citep{settles2009active},
core-set/$k$-center coverage~\citep{sener2018coreset}, uncertainty--diversity
hybrids~\citep{ash2020badge}, and farthest-point
sampling~\citep{eldar1997fps}. Whereas these select unlabeled
inputs, DISTIL selects failure \emph{modes} and adds a cross-round memory so
coverage accumulates over the whole run. Image-based failures are embedded
with a frozen R3M encoder~\citep{nair2022r3m}. Bridging demonstrations placed at
prescribed start states relate to reverse-curriculum reset
generation~\citep{florensa2017reversecurriculum} and
learning-to-reset~\citep{eysenbach2018leavenotrace}; DISTIL's start states
follow clustered failure evidence rather than a difficulty schedule.

\begin{figure*}[t]
\centering
\includegraphics[width=\textwidth]{figures/architecture.pdf}
\caption{\textbf{The DISTIL loop.} Rollouts of $f_\theta$ are monitored for
uncertainty via the per-step training loss; a vision LLM summarizes each
flagged failure's start, highest-loss ($t^\star$), and end frames, and a
reasoning LLM assigns root causes grounded in geometric facts (KAG);
failures are clustered into $k$ recurring modes with cross-round memory, the
LLM prescribes the round's corrective demonstration behind a feasibility
check with re-prescription, and the accepted demonstration augments the
dataset before the policy is updated.}
\label{fig:arch}
\end{figure*}

\section{Method}
\label{sec:method}

\subsection{Setup: Budgeted Interactive Imitation}

Each task is a finite-horizon decision process with states $s$, actions $a$
(discrete moves on GridWorld; relative joint-position deltas on the robot
tasks), and an expert $\pi^\star$ (a human on GridWorld; privileged
controllers on the robot tasks). The learner is \emph{any}
policy $f_\theta$ trained on an aggregated demonstration dataset
$\mathcal{D}$ by minimizing a per-step loss $\ell$,
\begin{equation}
\theta^\star=\arg\min_{\theta}\
\mathbb{E}_{(s,a)\sim\mathcal{D}}\big[\,\ell(f_\theta;s,a)\,\big],
\label{eq:il}
\end{equation}
where $\ell$ is the binary cross-entropy over the expert's action labels
for the GridWorld classifiers (pure CNN on images, MLP on states) and the
denoising ($v$-prediction) loss for the diffusion policies on the robot
tasks~\citep{ho2020ddpm,chi2023diffusionpolicy}.
Evaluated along a rollout at the learner's own executed actions, the loss
yields a per-step uncertainty signal
\begin{equation}
\ell_t=\ell\big(f_\theta;\,s_t,\,a^{\mathrm{nov}}_t\big),
\label{eq:signal}
\end{equation}
which scores how far the visited pair sits from the training
distribution~\citep{lee2025diffdagger}. Interactive learning proceeds in
rounds: round $r$ produces one corrective demonstration $d_r$ (a segment of
expert state--action pairs), and
\begin{equation}
\mathcal{D}_{r+1}=\mathcal{D}_r\cup d_r,
\qquad
\theta_{r+1}=\arg\min_{\theta}\
\mathbb{E}_{\mathcal{D}_{r+1}}[\,\ell\,].
\label{eq:loop}
\end{equation}
All methods receive the same hard budget of $B{=}20$ corrective
demonstrations, one per round; the only degree of freedom is how each
demonstration is chosen.

\subsection{Baseline Query Rules}
\label{sec:unified}

Every baseline instantiates one template: a scalar score is computed at each
visited state, and the expert takes over at the first threshold crossing,
\begin{equation}
\mathrm{Query}_m(s_t)=\mathbf{1}\big[\,\mathrm{score}_m(s_t)>\tau_m\,\big].
\label{eq:query}
\end{equation}
SafeDAgger scores the novice--expert action discrepancy
$\|a^{\mathrm{nov}}_t-a^{\mathrm{exp}}_t\|_2$ ($\tau{=}0.1$;
on GridWorld, the fraction of rollout steps whose action disagrees with the
expert's action labels)~\citep{zhang2017safedagger}. DropoutDAgger draws $N$ stochastic action
samples and fires when the fraction landing within a $\tau$-ball of the
expert action drops below $p$~\citep{menda2017dropoutdagger,gal2016dropout}.
EnsembleDAgger fires on the doubt (variance) of $M{=}5$ independently trained
members or on their mean
discrepancy~\citep{menda2019ensembledagger,lakshminarayanan2017ensembles}.
ThriftyDAgger fires on novelty or on learned task risk $1-Q_\psi(s,a)$, with
quantile-calibrated thresholds~\citep{hoque2021thriftydagger}. Stagger is a
uniform-random control that corrects one uniformly chosen recorded failure
per round. On the robot tasks we also run
Diff-DAgger~\citep{lee2025diffdagger}, whose score is the diffusion loss of
Eq.~\ref{eq:signal} itself: with $\hat F$ the empirical CDF of training-set
losses and $\alpha{=}0.99$,
\begin{equation}
\mathrm{Query}_{\mathrm{Diff}}(s_t)=
\mathbf{1}\big[\,\hat F(\ell_u)>\alpha\ \ \forall u\in(t{-}K,\,t]\,\big],
\label{eq:diffdagger}
\end{equation}
recalibrated at every retraining. The rosters differ for a structural
reason: Diff-DAgger's query rule \emph{is} the diffusion policy's own
training loss, so it requires a diffusion learner and does not apply to the
discrete-action CNN/MLP classifiers on GridWorld, where Stagger is the
matched uniform-random control; conversely, Stagger's random pick assumes
any visited state can be re-instantiated as a demonstration start, which
holds for grid layouts but not for contact-rich manipulation states, so on
the robot tasks Diff-DAgger, the loss-native rule, is the control (six
baselines overall, five per task). Whatever the
score, each rule decides \emph{when} to interrupt and corrects that rollout
from the crossing step; none compares failures, remembers corrected modes,
or chooses where the demonstration starts.

\subsection{The DISTIL Loop}
\label{sec:distil}

DISTIL replaces the per-state gate with a round-level analysis of the whole
failure set, run identically for every learner: from the pool of every
corrective demonstration the expert could give, it distills the few most
informative ones and spends the budget only on those.
Algorithm~\ref{alg:distil} gives the loop.

\paragraph{Uncertainty flagging and failure perception.}
Each round rolls the current policy and collects the failure set
$\mathcal{F}_r=\{f_1,\dots,f_N\}$. Along each episode the per-step loss
$\ell_t$ (Eq.~\ref{eq:signal}) is logged; steps whose loss exceeds an
adaptive threshold (the 95th percentile of the episode's loss distribution)
are flagged, and each failure retains its peak-loss step
\begin{equation}
t^\star_i=\arg\max_t \ell^{\,(i)}_t,
\qquad
\mathrm{peak}_i=\max_t \ell^{\,(i)}_t,
\label{eq:flag}
\end{equation}
taken as the failure point. A
vision language model then describes three frames per
failure (start, $t^\star_i$, and end) and produces a textual failure report
\citep{bai2025qwen3vl,liu2023reflect,duan2025aha}. A reasoning model
\citep{yang2025qwen3} assigns each failure a root cause (e.g., wrong
approach, no contact, wrong orientation, timeout) and a task phase. Both
stages are grounded by knowledge-augmented generation (KAG): a per-task
knowledge graph encoding the embodiment, object and goal geometry, workspace
bounds, a failure taxonomy, and per-failure-mode prescription rules is
rendered into the prompt.

\paragraph{Failure-mode clustering.}
Each failure is featurized by a descriptor anchored at $t^\star_i$. For
state-based robot runs this is the 6-D quaternion-free geometric vector
\begin{equation}
\phi_i=\big[\,p_{i,x},\ p_{i,y},\ \sin\theta_i,\ \cos\theta_i,\
\rho_i,\ \delta_i\,\big],
\label{eq:desc}
\end{equation}
with $(p_{i,x},p_{i,y},\theta_i)$ the object's planar pose,
$\rho_i=t^\star_i/T_i$ the task progress, and $\delta_i$ the end-effector
contact distance; for image-based runs the descriptor is a frozen R3M
embedding~\citep{nair2022r3m} of the same three frames, PCA-reduced to
$k'=\min(16,N{-}1)$ dimensions. On GridWorld, $\phi_i$ uses the agent's cell,
its signed offset to the goal, progress, and Manhattan distance-to-goal. The
standardized descriptors $\tilde X$ are clustered agglomeratively, with the
number of modes chosen by maximum mean silhouette,
\begin{equation}
k^\star=\arg\max_{k\in[2,\,k_{\max}]}\mathrm{sil}(k),
\label{eq:kstar}
\end{equation}
with $k_{\max}=\max\big(2,\min(6,N{-}1)\big)$.
Each cluster $C$ keeps its representative
$\mathrm{rep}(C)=\arg\min_{i\in C}\|\tilde X_i-\bar X_C\|_2$, centroid
geometry computed on raw poses, and mean peak loss $\bar L_C$; the dominant
cluster $C^\star$ is the largest, breaking ties by $\bar L_C$.

\paragraph{Cluster memory.}
Because the failure distribution shifts after every retraining, a persistent
mode could absorb the entire budget. DISTIL therefore remembers previously
corrected cluster centroids, penalizes recently covered regions with a
recency-discounted Gaussian kernel, and rotates the target among clusters
that rival the dominant one:
\begin{align}
P_{\mathrm{mem}}(c)&=
\sum_{(r_i,c_i)\in\mathrm{Mem}}
\gamma^{\,r-r_i}\,
\exp\!\Big(-\tfrac{\|c-c_i\|_2^2}{2\sigma^2}\Big),
\nonumber\\
C_{\mathrm{tgt}}&=\arg\max_{C:\,|C|\ge|C^\star|-1}
\big(\,\bar L_C-\lambda\,P_{\mathrm{mem}}(c_C)\,\big),
\label{eq:memory}
\end{align}
with $\gamma{=}0.6$, $\sigma{=}0.06$, $\lambda{=}1$. The prescription is
also conditioned on a small diverse context set $S$
($|S|\le\kappa{=}3$): the target representative is forced in, the
worst-peak-loss failure is seeded, and remaining slots are filled by
farthest-point selection in descriptor
space~\citep{sener2018coreset,eldar1997fps}.

\paragraph{Demonstration prescription.}
Given the clustered evidence, root causes, and cited context failures, the
LLM prescribes the round's single corrective demonstration together with a
self-reported confidence. The prescription is a hybrid decision with two
arms, both always in play. \emph{Targeted correction}: the LLM selects one
recorded failure; its scene is re-instantiated, the current policy replays to
its divergence point $t^\star$, and the expert takes over from there.
\emph{Bridging placement}: the LLM instead places the demonstration at
a new start state $\xi$ positioned between two or three cited failures of the
same mode, so one demonstration covers all of them. The bridge pose is
anchored to the target representative $A$: the commanded offset is
$L_2$-capped in position, clamped in orientation, and projected onto the KAG
workspace bounds $\mathcal{W}$,
\begin{align}
\Delta_{xy}&=\mathrm{cmd}.p_{xy}-p_{A,xy},
\quad
\Delta_{xy}\leftarrow
\Delta_{xy}\min\!\Big(1,\tfrac{\Delta_{\max}}{\|\Delta_{xy}\|_2}\Big),
\nonumber\\
\Delta_{\theta}&=\mathrm{clamp}\big(
\mathrm{wrap}_{\pi}(\mathrm{cmd}.\theta-\theta_A),\ \pm\theta_{\max}\big),
\nonumber\\
\xi&=\mathrm{clamp}_{\mathcal{W}}\big(
p_A+[\Delta_{xy},0],\ \ \theta_A+\Delta_{\theta},\ \ q_A\big),
\label{eq:prescribe}
\end{align}
with $(\Delta_{\max},\theta_{\max})=(0.06\,\mathrm{m},0.4\,\mathrm{rad})$ for
bridges on Push-T. On Lift the cube start is clamped to
$\Delta_{\max}=0.03$\,m in $x$ and $y$ with $\theta_{\max}=0$ (no yaw
randomization); on Door only the frame's $xy$ shifts, with
$(\Delta_{\max,x},\Delta_{\max,y})=(0.0135,0.013)$\,m and $\theta_{\max}=0$
(the bridge reuses the representative failure's orientation), and on Wipe
the randomized quantity is a whole dirt-marker path, so bridging is
infeasible and every prescription is a targeted correction. On GridWorld the
prescription is a corridor layout (start, goal, obstacles, and a cell
sequence) that the human expert is constrained to follow.

\paragraph{Feasibility check and re-prescription.}
Every prescription is validated before any expert effort is spent: robot
prescriptions must respect $\mathcal{W}$ and cite real recorded failures;
GridWorld corridors must stay on the grid, avoid the obstacles, and connect
start to goal under an $A^{\!\star}$/BFS path-validity check. An infeasible or empty prescription
is rejected and re-prescribed with the rejection fed back; after a bounded
number of attempts the round falls back deterministically to a targeted
correction of the nearest untried recorded failure, so a round always yields
a valid demonstration. The accepted demonstration
$d_r=\{(s_t,\pi^\star(s_t)):t\ge t^\star\}$ is added, the target centroid is
appended to the memory, and the policy is updated (Eq.~\ref{eq:loop}).

\begin{algorithm}[tb]
\caption{The DISTIL loop (one demonstration per round).}
\label{alg:distil}
\textbf{Input}: initial demos $\mathcal{D}_0$, expert $\pi^\star$, budget
$B{=}20$\\
\textbf{Output}: trained policy $f_\theta$
\begin{algorithmic}[1]
\STATE train $f_\theta$ on $\mathcal{D}_0$ (Eq.~\ref{eq:il});\ \
$\mathrm{Mem}\leftarrow\varnothing$
\FOR{$r=1$ \TO $B$}
  \STATE roll out $f_\theta$; flag uncertain episodes by $\ell_t$; record
  $t^\star_i,\mathrm{peak}_i$ (Eqs.~\ref{eq:signal},~\ref{eq:flag}); VLM
  describes start\,/\,$t^\star_i$\,/\,end frames; reasoning LLM assigns
  KAG-grounded root causes
  \STATE featurize (Eq.~\ref{eq:desc}); cluster into $k^\star$ modes
  (Eq.~\ref{eq:kstar}); rotate target via cluster memory
  (Eq.~\ref{eq:memory}); build context set $S$
  \STATE LLM prescribes the demonstration with confidence: targeted
  correction of a cited failure, or bridging start state $\xi$
  (Eq.~\ref{eq:prescribe}); re-prescribe while infeasible or empty; after
  bounded attempts, fall back to the nearest untried failure
  \STATE expert provides $d_r$;\ \
  $\mathcal{D}_r\leftarrow\mathcal{D}_{r-1}\cup d_r$;\ append target centroid
  to $\mathrm{Mem}$;\ update $f_\theta$ on $\mathcal{D}_r$ (Eq.~\ref{eq:loop})
\ENDFOR
\STATE \textbf{return} $f_\theta$
\end{algorithmic}
\end{algorithm}

\paragraph{Relation to the query rules.}
The baselines of Eqs.~\ref{eq:query}--\ref{eq:diffdagger} are the degenerate
round: a single-episode failure set, no clustering or memory, and the
prescription fixed to the targeted arm at the gate's own crossing step;
DISTIL spends the identical budget against the whole failure distribution
and the history of what was already corrected.

\paragraph{Bridging online and offline imitation.}
DAgger-family methods are \emph{online} interactive imitation learning: an
expert must be on call throughout training, because corrections are
collected while the policy rolls out. Classical behavior cloning is
\emph{offline}: collect demonstrations, then train. DISTIL bridges the two.
Its prescriptions are standing artifacts: the loop can run its rollout,
clustering, and prescription analysis on-policy and then pause; the
prescribed start configurations remain valid, and the expert can return days
later to record the prescribed demonstrations in one batch before the loop
resumes. Expert time is decoupled from policy training time: the expert
schedules a recording session instead of supervising every rollout.

\section{Experiments}
\label{sec:experiments}

We ask three questions. \textbf{(Q1)}~Under an identical 20-demonstration
budget, does DISTIL reach a higher final success rate than when-to-query
baselines? \textbf{(Q2)}~Does it collect more informative demonstrations,
and how does information gain relate to allocation across failure modes?
\textbf{(Q3)}~Are the LLM's prescription confidences predictive of realized
improvement, and are the discovered failure modes semantically meaningful?

\begin{table*}[t]
\centering
\small
\setlength{\tabcolsep}{4pt}
\caption{\textbf{Final held-out success rate (\%) at the 20-demonstration
budget}, mean\,$\pm$\,std over seeds (9 per cell on GridWorld, 5 on the
robot tasks). Columns abbreviate SafeDAgger, DropoutDAgger, EnsembleDAgger,
ThriftyDAgger; Stagger runs on GridWorld, Diff-DAgger on the robot tasks.
Best mean per row in \textbf{bold}: DISTIL is highest in all ten settings;
seed-level spreads overlap in the tighter cells (see text).}
\label{tab:main}
\begin{tabular}{llccccccc}
\toprule
Task & Obs.\ & Safe & Dropout & Ensemble & Thrifty & Stagger & Diff-DAgger & DISTIL (Ours) \\
\midrule
GridWorld $5{\times}5$ & state & $85.3 \pm 2.7$ & $84.9 \pm 2.5$ & $86.2 \pm 2.1$ & $86.8 \pm 2.0$ & $85.7 \pm 1.5$ & -- &$\mathbf{89.9 \pm 1.3}$ \\
GridWorld $5{\times}5$ & image & $86.1 \pm 2.8$ & $85.8 \pm 2.6$ & $85.7 \pm 2.2$ & $87.1 \pm 1.9$ & $86.6 \pm 2.3$ & -- &$\mathbf{89.6 \pm 1.8}$ \\
\addlinespace
Push-T & state & $82.0 \pm 6.8$ & $84.8 \pm 6.1$ & $85.9 \pm 5.8$ & $83.2 \pm 7.2$ & -- &$90.7 \pm 4.5$ & $\mathbf{96.1 \pm 4.5}$ \\
Push-T & image & $78.1 \pm 7.8$ & $82.1 \pm 6.9$ & $83.2 \pm 6.6$ & $79.3 \pm 8.1$ & -- &$89.0 \pm 4.8$ & $\mathbf{93.9 \pm 4.9}$ \\
\addlinespace
Lift & state & $99.2 \pm 1.6$ & $99.2 \pm 1.0$ & $99.2 \pm 1.0$ & $98.8 \pm 2.4$ & -- &$99.2 \pm 1.0$ & $\mathbf{100.0 \pm 0.0}$ \\
Lift & image & $99.6 \pm 0.8$ & $97.2 \pm 3.5$ & $98.8 \pm 1.6$ & $99.6 \pm 0.8$ & -- &$99.6 \pm 0.8$ & $\mathbf{100.0 \pm 0.0}$ \\
\addlinespace
Wipe & state & $88.0 \pm 2.5$ & $89.6 \pm 4.1$ & $90.8 \pm 4.3$ & $90.0 \pm 2.5$ & -- &$90.4 \pm 6.0$ & $\mathbf{95.5 \pm 6.0}$ \\
Wipe & image & $69.6 \pm 5.3$ & $83.2 \pm 6.8$ & $84.4 \pm 7.1$ & $69.2 \pm 9.0$ & -- &$89.6 \pm 3.2$ & $\mathbf{95.3 \pm 3.2}$ \\
\addlinespace
Door & state & $93.2 \pm 5.2$ & $92.8 \pm 2.7$ & $88.8 \pm 7.0$ & $89.6 \pm 3.9$ & -- &$95.2 \pm 4.3$ & $\mathbf{98.4 \pm 4.2}$ \\
Door & image & $92.4 \pm 3.2$ & $88.8 \pm 3.3$ & $86.0 \pm 10.9$ & $92.8 \pm 2.7$ & -- &$89.2 \pm 3.5$ & $\mathbf{99.2 \pm 3.4}$ \\
\bottomrule
\end{tabular}
\end{table*}

\begin{figure*}[t]
\centering
\includegraphics[width=0.92\textwidth]{figures/comparison_baselines.pdf}
\caption{\textbf{Success rate vs.\ demonstrations added} on GridWorld
$5{\times}5$ (image), Door (state), and Wipe (image); mean $\pm$ std over
seeds (9 for GridWorld, 5 for the robot tasks). Under the same 20-demo
budget, DISTIL dominates throughout training, with the largest margin on
image-based Wipe, where uncertainty-gated baselines plateau.}
\label{fig:comparison}
\end{figure*}

\subsection{Setup}

\paragraph{Tasks and experts.}
\textbf{GridWorld $5{\times}5$}: the agent moves from a start cell to a goal
cell with four discrete moves on a $5{\times}5$ grid containing three
obstacles; an episode succeeds when the agent reaches the goal. The
demonstrations are provided by a human expert; a shortest-path search is used
only as the feasibility check on prescribed layouts, never to produce
demonstrations. \textbf{Push-T}
(ManiSkill3~\citep{tao2024maniskill3,mu2021maniskill,gu2023maniskill2};
task from Implicit BC~\citep{florence2021implicitbc}): a Franka Panda with a
stick end-effector pushes a T-shaped block onto a fixed goal pose; the expert
is a privileged PPO policy. \textbf{Lift}, \textbf{Wipe}, and \textbf{Door}
(robosuite~\citep{zhu2020robosuite}, evaluation following
robomimic~\citep{mandlekar2021robomimic}): lift an object off the table, wipe
a mark off the table, and open a door, each with a scripted motion-planner
expert. Every task runs in a \emph{state} modality (privileged proprioception
and object state) and an \emph{image} modality (RGB observations hiding
object state).

\paragraph{Learners.}
On GridWorld the image learner is a pure CNN and the state learner an MLP,
both trained with cross-entropy on the expert's action labels. On all four
robot tasks every method shares one diffusion-policy
backbone~\citep{chi2023diffusionpolicy,ho2020ddpm}, a conditional 1-D U-Net
denoiser with a $v$-prediction objective, paired with an MLP state encoder or an
end-to-end finetuned R3M ResNet-18 image encoder with spatial-softmax
keypoints~\citep{nair2022r3m}. Within each cell all methods share the
identical learner and behavior-cloned initialization per seed; arms differ
only in how demonstrations are acquired.

\paragraph{Protocol and metrics.}
Every method receives a budget of exactly \textbf{20 expert demonstrations},
added one per round (Eq.~\ref{eq:loop}); the policy is retrained after each
addition and evaluated on a fixed held-out set shared across methods and
seeds. We run \textbf{9 seeds} per cell on GridWorld and \textbf{5 seeds} per
cell on the robot tasks. The primary metric is the final held-out success
rate (SR) at the budget. As analysis metrics we report the
\emph{per-demonstration information gain} (each round, before fine-tuning on
the newly collected demonstration, the current policy's per-step loss
evaluated on that demonstration) and, for DISTIL, the Pearson correlation
between prescription confidence and the realized change in success rate
after the demonstration is incorporated.

\begin{table}[t]
\centering
\small
\setlength{\tabcolsep}{2.4pt}
\caption{\textbf{Per-demonstration information gain} (pre-finetune policy
loss on each newly acquired demonstration; suboptimal demonstrations are
ruled out by construction, so higher = a demonstration from a region the
policy has not learned), mean\,$\pm$\,std pooled across a cell's seeds.
Abbreviations as in Table~\ref{tab:main}; D-DAg is Diff-DAgger. DISTIL is
the most informative in all ten settings.}
\label{tab:infogain}
\resizebox{\columnwidth}{!}{%
\begin{tabular}{llccccccc}
\toprule
Task & Obs. & Safe & Dropout & Ensemble & Thrifty & Stagger & D-DAg & DISTIL \\
\midrule
GridWorld & state & $2.55 \pm 1.68$ & $2.95 \pm 2.12$ & $1.33 \pm 1.18$ & $1.34 \pm 0.93$ & $1.83 \pm 1.11$ & -- &$\mathbf{3.55 \pm 2.51}$ \\
GridWorld & image & $2.46 \pm 1.61$ & $2.53 \pm 1.69$ & $1.57 \pm 1.55$ & $1.37 \pm 1.12$ & $1.88 \pm 1.43$ & -- &$\mathbf{3.21 \pm 2.33}$ \\
Push-T & state & $1.66 \pm 0.99$ & $2.36 \pm 1.63$ & $1.11 \pm 0.64$ & $1.10 \pm 0.65$ & -- &$1.57 \pm 1.10$ & $\mathbf{2.81 \pm 2.09}$ \\
Push-T & image & $2.04 \pm 1.11$ & $2.16 \pm 1.36$ & $1.06 \pm 0.64$ & $1.20 \pm 0.62$ & -- &$1.80 \pm 1.10$ & $\mathbf{2.82 \pm 1.72}$ \\
Lift & state & $2.23 \pm 1.36$ & $2.10 \pm 1.25$ & $1.12 \pm 0.74$ & $1.13 \pm 0.55$ & -- &$1.61 \pm 1.12$ & $\mathbf{2.64 \pm 1.65}$ \\
Lift & image & $2.18 \pm 1.40$ & $2.17 \pm 1.39$ & $1.00 \pm 0.55$ & $1.21 \pm 0.64$ & -- &$1.36 \pm 0.85$ & $\mathbf{2.93 \pm 1.67}$ \\
Wipe & state & $2.02 \pm 0.96$ & $2.38 \pm 1.55$ & $1.23 \pm 1.10$ & $1.18 \pm 0.81$ & -- &$1.43 \pm 0.80$ & $\mathbf{2.91 \pm 2.02}$ \\
Wipe & image & $2.50 \pm 1.47$ & $2.96 \pm 2.00$ & $1.43 \pm 0.89$ & $1.52 \pm 0.86$ & -- &$1.95 \pm 1.16$ & $\mathbf{3.62 \pm 2.19}$ \\
Door & state & $2.53 \pm 1.64$ & $3.10 \pm 2.21$ & $1.43 \pm 0.82$ & $1.42 \pm 0.90$ & -- &$1.84 \pm 1.11$ & $\mathbf{3.43 \pm 2.12}$ \\
Door & image & $2.35 \pm 1.40$ & $2.46 \pm 1.44$ & $1.24 \pm 0.76$ & $1.26 \pm 0.72$ & -- &$1.58 \pm 0.92$ & $\mathbf{3.00 \pm 1.98}$ \\
\bottomrule
\end{tabular}}
\end{table}

\begin{figure}[t]
\centering
\includegraphics[width=0.68\columnwidth]{figures/info_gain_boxplot.pdf}\\[2pt]
\includegraphics[width=0.58\columnwidth]{figures/confidence_vs_success.pdf}
\caption{\textbf{Top: per-demonstration information gain} (GridWorld
$5{\times}5$, image): the policy's pre-finetune loss on each newly acquired
demonstration; DISTIL's prescribed demonstrations carry the most new
information. \textbf{Bottom: prescription confidence predicts policy
improvement} (same setting): each point is one prescribed demonstration
($r=0.86$); across all ten task--modality settings, $r$ ranges from $0.82$
to $0.89$.}
\label{fig:infogain}
\label{fig:conf}
\end{figure}

\subsection{Main Results (Q1)}

Table~\ref{tab:main} reports final held-out success rates in all ten cells;
Figure~\ref{fig:comparison} shows learning curves on three representative
ones. \textbf{DISTIL attains the highest mean success rate
in every cell.} On state Push-T it reaches $96.1\%$ against $90.7\%$ for
Diff-DAgger, a $+5.4$-point margin on the task family Diff-DAgger was
designed for. Margins widen where observations are
hardest: on image Wipe, DISTIL reaches $95.3\%$ vs.\ $89.6\%$ for Diff-DAgger
and $69.6\%$ for SafeDAgger ($+25.7$ points), consistent with a gate that
keeps re-correcting the same contact-rich mode while the uncertainty-gated
curves of Figure~\ref{fig:comparison} flatten by mid-budget. On image Door,
DISTIL reaches $99.2\%$, $+6.4$ points over the best baseline, and on
GridWorld image it clears the best baseline (ThriftyDAgger) by $+2.5$ points
with the smallest variance in the column. Lift is near-saturated for all
methods, yet only DISTIL closes it completely ($100.0 \pm 0.0$ in both
modalities). With 5--9 seeds the seed-level spreads overlap the best
baseline in the tighter cells (GridWorld, Lift, state Door), so we do not
claim per-cell separation there; the evidence is the consistent
ranking (best mean in all ten cells) plus the wide margins on Push-T,
Wipe, and image Door.

\begin{figure}[t]
\centering
\includegraphics[width=0.62\columnwidth]{figures/clustering_modes_pusht.pdf}
\caption{\textbf{Failure modes discovered on image-based Push-T}: $k{=}3$
recurring modes (\emph{not-well-aligned}, \emph{no-contact}, and
\emph{badly-rotated}), with three representative frames each, annotated with
orientation error and contact distance. Coherent, well-separated clusters let
the LLM prescribe one targeted demonstration per mode.}
\label{fig:clusters}
\end{figure}

\subsection{Information Gain and Its Allocation (Q2)}

Table~\ref{tab:infogain} and Figure~\ref{fig:infogain} quantify how much new
information each acquired demonstration carries: each round, before
fine-tuning on the newly collected demonstration, we evaluate the current
policy's per-step loss on it. A high pre-finetune loss can mean only two
things: (a)~the demonstration covers a region the current policy has not
learned (diverse data from an underrepresented region), or (b)~the
demonstration is suboptimal or invalid. Case~(b) is ruled out by
construction: every prescription passes the feasibility check before any
demonstration is collected, and the demonstrations themselves are produced by
the expert. High pre-finetune loss therefore identifies demonstrations from
underrepresented regions, and the per-demonstration information gain measures
exactly the diversity a method's chosen demonstrations add. DISTIL's
prescribed demonstrations have the highest mean information gain in all ten
cells (e.g., $3.62$ vs.\ next-best $2.96$ on image Wipe; $3.55$ vs.\ $2.95$
on state GridWorld), with the widest upper spread in the boxplot: the
prescription deliberately seeks out what the policy cannot yet do. One
caveat: SafeDAgger and DropoutDAgger collect higher raw
information gain than EnsembleDAgger and ThriftyDAgger, yet the
information-gain orderings do not match the success-rate orderings of
Table~\ref{tab:main}, so high per-demonstration gain is necessary but not
sufficient. What converts it into success rate is its \emph{allocation
across failure modes}: a gate that repeatedly harvests high-loss
demonstrations from the \emph{same} mode accumulates redundant information,
whereas DISTIL's clustering and cluster memory spread high-loss
demonstrations over distinct, not-yet-corrected modes; this reading
matches the flattened baseline curves of
Figure~\ref{fig:comparison}.

\subsection{Confidence Calibration and Discovered Modes (Q3)}

Each prescription carries a self-reported confidence. Why should the
LLM's prescription be trusted in each round? The success-rate feedback
arrives late: only after the demonstration is collected, the policy
retrained, and new rollouts evaluated. The prescription confidence is
reported at prescription time, blind to that future outcome and to the
success-rate trajectory. Figure~\ref{fig:conf} shows that this blind,
self-reported confidence correlates with the realized policy improvement on
image GridWorld ($r=0.86$); the same computation per cell yields $r=0.88$
(state GridWorld), $0.87/0.88$ (state/image Push-T), $0.88/0.89$ (Lift),
$0.82/0.86$ (Wipe), and $0.83/0.82$ (Door), a tight $0.82$--$0.89$ range
across all ten settings. A confident prescription is therefore evidence,
available immediately, that the round will help. Low-confidence
prescriptions ($<60\%$) yield near-zero improvement while high-confidence
ones ($\ge 90\%$) yield the largest gains, so a practitioner could audit or
veto low-confidence prescriptions before expert effort is committed.
Figure~\ref{fig:clusters} shows the modes discovered on image-based Push-T:
three coherent, well-separated clusters, internally consistent in
orientation error and contact distance. These modes are the units over
which the budget is allocated.

\section{Conclusion}
\label{sec:conclusion}

Interactive imitation learning has long asked \emph{when} to query the
expert. Under a fixed budget that is the wrong question: performance is set
by \emph{which} failures are corrected and \emph{how} each demonstration is
placed. DISTIL operationalizes this as demonstration distillation: flag
uncertain rollouts with the learner's own loss, root-cause failures with
KAG-grounded foundation models, cluster them into modes with a cross-round
memory, and prescribe each demonstration through a feasibility-checked LLM
loop that corrects a chosen failure in place or bridges several related
ones. Across five tasks, two
observation modalities, and four learner architectures, the same loop attains
the best mean success rate against six interactive baselines in every setting
at a 20-demonstration budget, and its demonstrations carry the most new
information, spread by design across distinct failure modes.

\paragraph{Limitations.}
On the robot tasks our experts are oracles (a privileged PPO policy and
scripted motion planners); noisy or inconsistent human demonstrators would
add variance this expert model does not
capture~\citep{mandlekar2020iwr,liu2023sirius}. Bridging placements require
a resettable simulator, and all results are in simulation. The loop inherits
the quality of its foundation models: weaker VLM or LLM outputs degrade
root-causing and prescription; KAG grounding, workspace clamps, and the
feasibility loop bound, but do not eliminate, the damage of a bad
prescription. Each round also pays the
wall-clock cost of the reasoning pipeline (frame description, root-cause
analysis, clustering, prescription, and possible re-prescription), an
inference cost the per-state gates do not incur, and each new task family
requires authoring its knowledge graph before the loop can run. The loop
commits to a single demonstration per round and cannot exploit an expert
willing to record several complementary corrections at once, and the
descriptor $\phi_i$ and memory constants are designed, not learned. Letting
the VLM's failure taxonomy define the clustering space and extending
prescription to human-gated hardware
settings~\citep{kelly2019hgdagger} are natural next steps.

\bibliographystyle{aaai2027}
\bibliography{references}

\end{document}
